\chapter{Fundamentals and Related Work}
\label{chap:fundamentals}

This chapter discusses the concepts and mathematical background needed to
understand point cloud upsampling and its use in three-dimensional
perception. It begins with point representation and sampling, then explains
classification and detection before introducing the upsampling methods.
The final section defines geometric and task metrics. The purpose is to
explain how the methods work and how their outputs should be interpreted;
the particular datasets, training settings and evaluation configurations
are introduced in Chapter~\ref{chap:setup}.

\section{Point Clouds and 3D Perception}
\label{sec:point_clouds}

Point clouds provide spatial observations; classification and detection
assign semantic meaning to those observations. In object classification,
a classifier receives the points of one segmented object and predicts one
category~\cite{qi2017pointnet}. In three-dimensional object detection, a
detector receives a scene and predicts object instances, their categories,
confidence scores and oriented bounding boxes
\cite{shi2019pointrcnn,yin2021centerpoint}. A detector must therefore
localise objects as well as recognise them. The distinction connects the
point representation introduced below to the classification networks in
Section~\ref{sec:classification} and the detectors in
Section~\ref{sec:detection}. Upsampling supplies denser coordinates to
either task; it does not itself assign an object label or a detection box
\cite{yu2018punet,qian2021pugcn}.

\subsection{Definition and Properties of Point Clouds}
\label{subsec:pc_definition}

A point cloud is an unordered finite set of discretely sampled points in
three-dimensional space~\cite{qi2017pointnet}. When only geometric
coordinates $x_i$, $y_i$ and $z_i$ are used, a point cloud consisting
of $N$ points is written as
\begin{equation}
  \mathcal P=\{\mathbf p_i\}_{i=1}^N,
  \qquad \mathbf p_i=(x_i,y_i,z_i)^{\mathsf T}\in\mathbb R^3.
  \label{eq:pointcloud}
\end{equation}
Here $N$ is the number of stored points, $i$ indexes a point, and
$(x_i,y_i,z_i)$ are its Cartesian coordinates. The superscript
$\mathsf T$ denotes transpose, so that $\mathbf p_i$ is a column vector;
$\mathbb R^3$ is three-dimensional real coordinate space. A point gives a
location, not a small solid object. Connections, faces and a surface normal
are not supplied by the coordinate set itself. A mesh, in contrast,
explicitly records connections between vertices~\cite{qi2017pointnet}.

Figure~\ref{fig:random-point-cloud} illustrates this definition with a
random point set. The example has no object category and no experimental
meaning: it shows only positions in a coordinate system. The highlighted
point has three coordinates just like every other point. The apparent
distance between points in a projected image must not be confused with
their full three-dimensional distance.
\thesisfigure{random_point_cloud}{An unordered point cloud}{A synthetic
random set of 192 points in the unit cube, used only to illustrate
three-dimensional coordinates and irregular point spacing. All points are
shown; the highlighted point is one member of the same set. Coordinates
are dimensionless. This is a conceptual example, not an object sample or
an experimental result.}{fig:random-point-cloud}

If the sensor additionally provides reflectance, colour or normal vectors,
the spatial coordinates can be stored together with an attribute vector as
\begin{equation}
  \mathbf u_i=[\mathbf p_i^{\mathsf T},\mathbf a_i^{\mathsf T}]^{\mathsf T}
  \in\mathbb R^{3+d_a}.
  \label{eq:point_attributes}
\end{equation}
Here $\mathbf a_i\in\mathbb R^{d_a}$ contains $d_a$ attributes, and
$\mathbf u_i$ is the extended record. This notation keeps the spatial
point $\mathbf p_i\in\mathbb R^3$ separate from its features. An XYZI record,
for example, contains $(x_i,y_i,z_i,\iota_i)$ with measured reflectance
$\iota_i$~\cite{geiger2013kitti}. Distances between positions use the first
three coordinates; adding reflectance to a Euclidean distance without a
defined scaling would mix different quantities.

Point clouds possess, first of all, the property of being
\emph{unordered}. Writing the records as rows of a matrix
$\mathbf U\in\mathbb R^{N\times(3+d_a)}$ is a storage convention and
does not express a natural ordering. For a set-level output such as a
classification score, a function $f$ must satisfy
\begin{equation}
 f(\{\mathbf p_1,\ldots,\mathbf p_N\})
 =f(\{\mathbf p_{\pi(1)},\ldots,\mathbf p_{\pi(N)}\})
 \label{eq:permutation_invariance}
\end{equation}
for every permutation $\pi$ of the indices $1,\ldots,N$. This is
permutation invariance. For a per-point output, such as foreground
segmentation, reordering the inputs should reorder the outputs in the same
way; this is permutation equivariance. A shared pointwise mapping followed
by symmetric aggregation implements the invariant case in
PointNet~\cite{qi2017pointnet}. Section~\ref{subsec:pointnet} explains how the network uses this construction.

Point clouds often exhibit \emph{irregular sampling}. The representation
does not require a fixed image-like lattice, and local distances and
neighbourhood sizes can vary substantially~\cite{qi2017pointnetpp}.
Measured point clouds are also often \emph{incomplete}, because a sensor
records visible surfaces subject to occlusion and acquisition limitations
\cite{geiger2013kitti}. Unordered storage, irregular density and incomplete
coverage are different properties. A randomly reordered cloud preserves
its geometry; removing points can change both its sampling density and its
coverage. These distinctions lead to two different observation processes.

An object-level cloud represents one segmented shape; surface sampling
from a mesh can cover regions that would be hidden from a single viewpoint.
A scene-level scan instead contains returns from several objects, ground
and background. Surface coverage depends on the observation process, not
on the definition of a point cloud
\cite{wu2015modelnet,geiger2013kitti,qi2017pointnet}.
The concrete object dataset, normalisation and scene data preparation are
specified in Sections~\ref{sec:fc-4-3-1} and~\ref{sec:fc-4-3-2}.

\subsection{Differences Between Object-Level and Scene-Level Point Clouds}
\label{subsec:domain_gap}

Table~\ref{tab:object_vs_scene} compares the two
observation processes. It describes common surface-sampling and sensor
settings; it does not specify a particular experimental split or point
count~\cite{geiger2013kitti,shi2019pointrcnn,wu2015modelnet,qi2017pointnet}.
\begin{table}[!t]
 \centering
 \caption[Object-level point clouds versus LiDAR scenes]{Structural
 differences between object-level surface samples and LiDAR scenes.
 Surface completeness and normalisation describe common preparation
 conventions, not requirements of either representation
 \cite{geiger2013kitti,shi2019pointrcnn,wu2015modelnet,qi2017pointnet}.}
 \label{tab:object_vs_scene}
 \begin{tabular}{L{0.24\textwidth}L{0.32\textwidth}L{0.32\textwidth}}
 \toprule
 \textbf{Comparison dimension}&\textbf{Object-level point cloud}&\textbf{LiDAR scene}\\
 \midrule
 Data source&Sampling of a CAD mesh surface&Real sensor returns\\
 Spatial content&A single, already segmented object&Multiple objects, ground and background\\
 Surface completeness&Can cover the complete available surface&Covers sensor-visible surfaces\\
 Coordinate scale&Often centred and unit-normalised&Real metric coordinates\\
 Sampling regularity&Controlled surface sampling&Determined by scan angles, range and occlusion\\
 Primary task&Output of one object category&Output of multiple categories and 3D boxes\\
 Dominant error consequence&Altered class-discriminative features&Altered class, centre, size and orientation\\
 \bottomrule
 \end{tabular}
\end{table}

In the last row, \emph{error} refers to a distortion of the input geometry,
such as missing surface support, displaced points or spurious added points.
Its \emph{consequence} is a possible change in the features used for a task.
For classification, this can alter the evidence distinguishing categories.
For detection, it can additionally affect the estimated centre, dimensions
or orientation of an object~\cite{qi2017pointnet,shi2019pointrcnn}.
The row describes possible pathways from geometric distortion to a task
error; it does not assert that every upsampling operation causes them.

The table therefore explains why the same geometric operation need not have
the same task effect in both settings. A category label summarises one input shape;
a detector must additionally separate instances and localise them within a
scene~\cite{shi2019pointrcnn,qi2017pointnet}. To understand where additional
points enter those decisions, the next two sections follow the processing
chain from coordinates to class scores and then from scene features to
three-dimensional boxes.

\FloatBarrier
\section{Point Cloud Classification}
\label{sec:classification}

\subsection{Problem Formulation}
\label{subsec:cls_formulation}

The input of point cloud classification is the point set $\mathcal{P}$ of a
single object and the output is its category $y$~\cite{qi2017pointnetpp,qi2017pointnet}. Here $y\in\{1,\ldots,C\}$ and $C$ is the number of classes. The subscript $\theta$ denotes learned parameters. A classification network
$f_\theta$ produces the logits of $C$ categories,
%
\begin{equation}
    \label{eq:logits}
    \mathbf{z} = f_\theta(\mathcal{P}) \in \mathbb{R}^{C}.
\end{equation}
%
Here $\mathbf z$ is a vector of real-valued class scores, or logits. After softmax, the predicted probability $\hat p_c$ of class $c$ is
%
\begin{equation}
    \label{eq:softmax}
    \hat{p}_c = \frac{\exp(z_c)}{\sum_{j=1}^{C} \exp(z_j)} .
\end{equation}
%
Here $z_c$ is the score of class $c$, $j$ indexes all candidate classes, and $\exp$ is the exponential function. The denominator makes the probabilities sum to one. Given the ground-truth label $y$, the cross-entropy objective used for classification~\cite{qi2017pointnet} is
%
\begin{equation}
    \label{eq:cross_entropy}
    \mathcal{L}_{\mathrm{cls}} = -\log \hat{p}_y ,
\end{equation}
%
where $\log$ is the natural logarithm and $\hat p_y$ is the probability assigned to the true class. A confident correct prediction has a small loss; assigning little probability to the true class gives a large loss. For $M$ training samples the empirical risk is
%
\begin{equation}
    \label{eq:empirical_risk}
    \mathcal{L}_{\mathrm{train}}
    = \frac{1}{M} \sum_{m=1}^{M} \mathcal{L}_{\mathrm{cls}}^{(m)} .
\end{equation}
%
The superscript $(m)$ indexes a sample and $\mathcal L_{\mathrm{train}}$ is its sample-averaged training objective. These equations express the classification objective of PointNet-family networks~\cite{qi2017pointnetpp,qi2017pointnet}. The classification decision is $\hat{y} = \arg\max_c \hat{p}_c$. This
formulation shows that upsampling yields a benefit only if it changes the
features extracted by the network in such a way that the score of the true
class is raised relative to the scores of the remaining classes. Defining
the classification margin as
%
\begin{equation}
    \label{eq:margin}
    \Delta_{\mathrm{cls}}(\mathcal{P}) = z_y - \max_{c \neq y} z_c ,
\end{equation}
%
a positive margin means that the true class has the largest score. Here
$z_y$ is the true-class logit and the maximum selects the strongest competing
class; a zero margin is a tie. For a fixed classifier, a change in point
count alone imposes no constraint on this margin. The margin can increase or decrease when the input changes. A change
across zero can correct a previous classification error or introduce a
new one. These statements follow from the decision rule; they are not
additional measured results. The
following architectures explain how the point coordinates become the scores
used by that rule.

\subsection{PointNet}
\label{subsec:pointnet}

PointNet processes an unordered point set directly through a shared
pointwise mapping, symmetric aggregation and a classification head
\cite{qi2017pointnet}. The shared multilayer perceptron (MLP) applies
the same learned transformation to every point. Its output is a feature
vector describing that point in a learned space. A channelwise max
operation then retains the largest response in each feature channel
across the whole set. Because the maximum is unchanged by reordering,
the resulting global feature satisfies the permutation invariance in
Equation~\eqref{eq:permutation_invariance}. The classification head maps
that feature to category scores.

PointNet also uses transformation networks to align input coordinates
and intermediate features in learned spaces. The feature transformation
is regularised towards an orthogonal matrix, which discourages severe
distortion without imposing an exact rigid rotation
\cite{qi2017pointnet}. These transformations complement the shared
mapping and pooling; they do not introduce an ordering of the points.

The global maximum is determined by the points attaining the largest
response in each channel. Qi et al. describe this property through a
critical-point set~\cite{qi2017pointnet}. Repeating an identical
pointwise feature does not change a channel's maximum. However, the
shared mapping before pooling does not explicitly aggregate the relative
positions of neighbouring points. PointNet++ addresses this limitation
by applying local PointNet operations within nested metric regions
\cite{qi2017pointnetpp}. This progression from global aggregation to
local geometric aggregation connects the two architectures.

\subsection{PointNet++}
\label{subsec:pointnetpp}

PointNet++~\cite{qi2017pointnetpp} learns local-to-global features
progressively by means of hierarchical sampling, radius grouping and local
PointNet aggregation. A set abstraction layer usually comprises the three
steps of sampling, grouping and feature extraction.

Figure~\ref{fig:pointnet-principle} connects these operations in the feature hierarchy.
\thesisfigure{pointnet_principle}{PointNet++ principle}{Sampling, local
grouping and shared feature aggregation in PointNet++, redrawn from the
published architecture~\cite{qi2017pointnetpp}.}{fig:pointnet-principle}

The first step selects $M < N$ centroids by
\ac{FPS}, as used in PointNet++~\cite{qi2017pointnetpp}. If the set of already selected centroids is
$\mathcal{M}_t$, the next centroid is
%
\begin{equation}
    \label{eq:fps}
    \mathbf{c}_{t+1}
    = \arg\max_{\mathbf{p}_i \in \mathcal{P}}\;
      \min_{\mathbf{c} \in \mathcal{M}_t} \big\| \mathbf{p}_i - \mathbf{c} \big\|_2 .
\end{equation}
%
Here $t$ counts selection steps, and $\mathcal M_t$ contains the
centroids already selected. The inner minimum gives a candidate point's
distance to its nearest selected centroid. The outer maximum chooses the
candidate with the largest such distance, namely the least-covered
remaining point. Starting from an initial centroid, the procedure repeats
until $M$ centroids have been selected~\cite{qi2017pointnetpp}.
An initial-point rule and a rule for equal distances specify the selection
completely. The operation selects existing points; it does not create
new coordinates.

The second step constructs a local neighbourhood
around each centroid $\mathbf{c}_j$. The radius grouping takes the form
%
\begin{equation}
    \label{eq:ball_query}
    \mathcal{N}_j(R) = \big\{ \mathbf{p}_i : \| \mathbf{p}_i - \mathbf{c}_j \|_2 \le R \big\}.
\end{equation}
%
Here $R>0$ is a radius in the coordinate units of the input; $j$ indexes centroids. A $k$-nearest-neighbour query instead fixes the number $k$ of neighbours and permits the neighbourhood radius to vary. The two alternatives are discussed in PointNet++~\cite{qi2017pointnetpp}. In order to
obtain a translation-invariant local description, relative coordinates are
obtained by subtracting the centroid,
%
\begin{equation}
    \label{eq:relative_coords}
    \Delta \mathbf{p}_{ij} = \mathbf{p}_i - \mathbf{c}_j .
\end{equation}
%
Unlike a $k$ nearest neighbour query, the ball query of
\eqref{eq:ball_query} guarantees a fixed metric scale, so that the learned
features carry a consistent spatial scale in regions of differing
density. For normalised coordinates this scale is relative to the normalisation radius, not a distance in metres.

The third step applies a shared mapping and
a symmetric pooling within each local set, expressing the grouped points
relative to their centroid and producing one feature vector per centroid according to
%
\begin{equation}
    \label{eq:sa_layer}
    \mathbf{f}'_j
    = \max_{\mathbf{p}_i \in \mathcal{N}_j}
      \phi_\theta\big( [\, \Delta \mathbf{p}_{ij},\, \mathbf{f}_i \,] \big),
\end{equation}
%
where $\phi_\theta$ is a shared MLP with learned parameters $\theta$,
$\mathbf f^{\prime}_j$ is the output feature, $[\,\cdot,\cdot\,]$ denotes concatenation, $\mathcal N_j$ is the ball-query set, and $\mathbf{f}_i$ is the point feature of the preceding layer; the first
layer may use the coordinates alone or the coordinates together with the
normal vectors. Stacking several set abstraction layers steadily reduces the
number of centroids while enlarging the receptive field, and finally yields
an object-level feature. In order to accommodate a non-uniform density,
PointNet++ can employ \ac{MSG}. Given several radii $R_1,\dots,R_L$, the
local features are extracted separately and concatenated as
%
\begin{equation}
    \label{eq:msg}
    \mathbf{f}^{\mathrm{MSG}}_j
    = \mathop{\Big\Vert}\limits_{\ell=1}^{L}\;
      \max_{\mathbf{p}_i \in \mathcal{N}_j(R_\ell)}
      \phi_{\theta_\ell}\big( [\, \mathbf{p}_i - \mathbf{c}_j,\, \mathbf{f}_i \,] \big).
\end{equation}
%
Here $L$ is the number of scales, $\ell$ indexes a scale, $R_\ell$ is its radius, and $\theta_\ell$ denotes the parameters of that branch; $\Vert$ denotes feature concatenation. Single-scale grouping (SSG) uses one radius per layer, whereas multi-scale grouping (MSG) concatenates several radii~\cite{qi2017pointnetpp}. A small radius attends to
detail, whereas a large radius provides a wider context whenever the number
of local points is insufficient.

\FloatBarrier
\section{3D Object Detection from Point Clouds}
\label{sec:detection}

The preceding classifier maps one point cloud to one category. Detection
extends this task to a scene containing several objects and background:
it must decide which points support an object, distinguish instances, and
estimate each object's position and extent~\cite{shi2019pointrcnn}.
An oriented three-dimensional box can be written as
\begin{equation}
 \mathbf b=(x_c,y_c,z_c,l,w,h,\psi)^{\mathsf T},
 \label{eq:box}
\end{equation}
where $(x_c,y_c,z_c)$ is the centre, $(l,w,h)$ are length, width and height,
and $\psi$ is yaw about the vertical axis. A detector returns
\begin{equation}
 \hat{\mathcal B}=\{(\hat{\mathbf b}_j,\hat c_j,\hat s_j)\}_{j=1}^{n_\mathrm{det}}.
 \label{eq:detections}
\end{equation}
Here $n_\mathrm{det}$ is the number of retained detections, $j$ indexes one
detection, $\hat c_j$ is its predicted class and $\hat s_j$ its confidence.
The box notation is independent of a particular coordinate convention;
the choice of vertical axis must be stated when writing its regression
targets~\cite{shi2019pointrcnn}.

Both detection families considered below learn features, classify candidate
objects and regress box parameters. They differ in the representation from
which candidates are formed. PointRCNN predicts proposals from foreground
points; CenterPoint predicts centre peaks from a spatial feature map
\cite{shi2019pointrcnn,yin2021centerpoint}. Figure~\ref{fig:detector-principles}
shows where these processing paths diverge. The distinction matters for
upsampling because adding points changes the input to a sequence of
sampling and aggregation operations, not directly the final boxes.
\thesisfigure{detector_principles}{Point-based and centre-based detection}{
Processing paths of PointRCNN and a voxel-based CenterPoint configuration,
redrawn from the method descriptions~\cite{shi2019pointrcnn,yin2021centerpoint}.
The diagram shows the principal representations; it omits layer sizes and
implementation-specific input limits.}{fig:detector-principles}

\subsection{PointRCNN}
\label{subsec:pointrcnn}

PointRCNN is a two-stage, point-based detector. Its first stage uses a
PointNet++ backbone to obtain pointwise semantic features, segments points
into foreground and background, and generates proposals directly from
foreground points. Its second stage pools the points around each proposal
and refines that proposal in local coordinates~\cite{shi2019pointrcnn}.
This connects the local feature hierarchy of
Section~\ref{subsec:pointnetpp} to instance localisation.

Let $\mathbf f_i\in\mathbb R^F$ be the $F$-dimensional feature of point
$\mathbf p_i$. A learned scalar mapping $g_\theta$ produces a foreground
probability
\begin{equation}
 \hat q_i=\operatorname{sigmoid}(g_\theta(\mathbf f_i)),
 \qquad\operatorname{sigmoid}(u)=\frac{1}{1+e^{-u}}.
 \label{eq:fg_prob}
\end{equation}
PointRCNN derives foreground labels from annotated three-dimensional boxes
and uses focal loss to address the foreground--background imbalance
\cite{shi2019pointrcnn,lin2017focal}. For binary label $q_i\in\{0,1\}$,
define the probability of the correct binary class and its weighting as
\begin{equation}
 \begin{aligned}
 p_{t,i}&=q_i\hat q_i+(1-q_i)(1-\hat q_i),\\
 \alpha_{t,i}&=q_i\alpha_f+(1-q_i)(1-\alpha_f).
 \end{aligned}
 \label{eq:focal_targets}
\end{equation}
The per-point focal loss is
\begin{equation}
 \ell_{\mathrm{seg},i}=-\alpha_{t,i}(1-p_{t,i})^{\gamma_f}\log p_{t,i}.
 \label{eq:focal}
\end{equation}
Here $\alpha_f\in(0,1)$ controls class weighting and $\gamma_f\ge0$
controls suppression of easily classified examples. If the correct class
already has probability close to one, the factor
$(1-p_{t,i})^{\gamma_f}$ reduces its contribution. For $q_i=0$ the formula
uses $1-\hat q_i$, so it supervises background as well as foreground.
This is the binary focal-loss formulation of Lin et al.; it does not
specify the hyperparameters of the experiments in this
thesis~\cite{lin2017focal}.

A foreground point provides a local reference from which PointRCNN
predicts the offset to the object centre, box dimensions and orientation.
For horizontal position and yaw, PointRCNN combines a discrete bin
classification with a continuous residual within the selected bin
\cite{shi2019pointrcnn}. To express the centre-offset construction without
mixing camera and LiDAR axes, let $\zeta$ denote either horizontal offset.
For a horizontal offset $\zeta\in[-S,S)$ and bin width $\delta>0$,
\begin{equation}
 b^\star=\left\lfloor\frac{\zeta+S}{\delta}\right\rfloor,
 \quad c_b=-S+(b+\tfrac12)\delta,
 \quad r^\star=\frac{\zeta-c_{b^\star}}{\delta}.
 \label{eq:bin_decomposition}
\end{equation}
Here $S>0$ is the half-width of the search range, $b^\star$ the target bin,
$c_b$ the centre of bin $b$, and $r^\star$ its dimensionless residual.
The floor operation selects a discrete interval. With predicted bin
$\hat b$ and predicted residual $\hat r_{\hat b}$, the offset is recovered
as $\hat\zeta=c_{\hat b}+\delta\hat r_{\hat b}$. Angular regression uses
the analogous construction over an angular interval. The original paper
uses camera-frame $x$ and $z$ as horizontal axes; this does not make
$z$ horizontal in a LiDAR coordinate convention~\cite{shi2019pointrcnn}.

The bin label is supervised by cross-entropy and the residual by smooth
$L_1$ regression. For scalar residual error $e$, the latter is
\begin{equation}
 \operatorname{smooth}_{L_1}(e)=
 \begin{cases}\tfrac12e^2,&|e|<1,\\|e|-\tfrac12,&|e|\ge1.\end{cases}
 \label{eq:smooth_l1}
\end{equation}
It is quadratic near zero and linear for larger errors. Thus the network
first learns a coarse location and then a finer correction, while other
box terms such as dimensions are regressed separately
\cite{shi2019pointrcnn}. Proposal filtering removes redundant candidates
before the refinement stage.

The second stage pools point coordinates and features within an enlarged
proposal region. For proposal centre $\mathbf c$ and yaw $\psi$, a
$z$-up notation for its local coordinates is
\begin{equation}
 \mathbf p_i^\mathrm{can}=\mathbf R_z(-\psi)(\mathbf p_i-\mathbf c),
 \qquad
 \mathbf R_z(\varphi)=
 \begin{bmatrix}
 \cos\varphi&-\sin\varphi&0\\
 \sin\varphi&\cos\varphi&0\\
 0&0&1
 \end{bmatrix}.
 \label{eq:canonical}
\end{equation}
Here $\mathbf R_z(\varphi)$ rotates through angle $\varphi$ around the
vertical axis. Translation removes the proposal centre and inverse yaw
aligns its heading. The refinement network then predicts residual box
parameters and confidence from the aligned region. PointRCNN uses the
corresponding transformation about the vertical axis of its working frame
\cite{shi2019pointrcnn}.

These stages explain two possible routes by which generated points can
affect a point-based detector. They can change the pointwise features and
foreground proposals, and they can change the points pooled for refinement.
Moreover, if an implementation samples an input set $\mathcal Y$ to a
budget $B_p\le|\mathcal Y|$, it actually processes
\begin{equation}
 \mathcal P_\mathrm{det}=\mathcal S_{B_p}(\mathcal Y),
 \qquad |\mathcal P_\mathrm{det}|=B_p,
 \label{eq:point_budget}
\end{equation}
where $\mathcal S_{B_p}$ denotes the chosen sampling operator. PointRCNN
uses input point sampling~\cite{shi2019pointrcnn}; the consequence is that
file-level cardinality need not equal the number of points consumed.
The implemented full-frame sampling and the separate region-split
controls are distinguished in Sections~\ref{sec:fc-4-8-1}
and~\ref{sec:fc-4-12-3}. The next detector aggregates points into spatial cells before
forming object candidates, introducing a different kind of reduction.

\subsection{CenterPoint}
\label{subsec:centerpoint}

CenterPoint extends the centre-keypoint idea of CenterNet to
three-dimensional detection~\cite{yin2021centerpoint,zhou2019objectsaspoints}.
It predicts a class-specific centre heatmap in bird's-eye view and regresses
box attributes at the predicted centres. The centre head can operate with
voxel or pillar encoders; voxelisation is a backbone choice rather than
the definition of every CenterPoint variant~\cite{yin2021centerpoint}.
The voxel-based path is explained here to show how point density and spatial
occupancy differ.

With voxel dimensions $\mathbf v=(v_x,v_y,v_z)^{\mathsf T}$, where each
entry is positive, and a lower spatial bound $\mathbf o$, a point is
assigned to the grid cell
\begin{equation}
 \mathbf k_i=\left\lfloor(\mathbf p_i-\mathbf o)\oslash\mathbf v\right\rfloor.
 \label{eq:voxel_index}
\end{equation}
Here $\oslash$ denotes componentwise division and the floor acts on each
coordinate. The integer vector $\mathbf k_i$ identifies a voxel. A voxel
feature encoder combines the points within each occupied cell, and sparse
convolutions propagate information across the occupied grid
\cite{zhou2018voxelnet,yan2018second}.

For example, the learned voxel encoding used by VoxelNet combines each
point's features with its displacement from the voxel centroid before
symmetric pooling~\cite{zhou2018voxelnet}. Abstracting this operation,
\begin{equation}
 \mathbf f_k=\max_{i:\mathbf k_i=k}
 \phi_\theta([\mathbf u_i,\mathbf p_i-\bar{\mathbf p}_k]).
 \label{eq:vfe_general}
\end{equation}
The vector $\bar{\mathbf p}_k$ is the mean position in voxel $k$,
$\mathbf u_i$ is the coordinate-and-attribute record from
Equation~\eqref{eq:point_attributes}, and $\phi_\theta$ is a shared learned
mapping. The maximum operates channelwise. Mean encoders instead average
retained point records; these are alternative implementations, not the
same operation as VoxelNet's learned encoder. The spatial backbone converts
the encoded grid into the bird's-eye-view feature map used by the centre
head~\cite{yin2021centerpoint}.

For one annotated centre projected onto a feature grid at integer position
$\boldsymbol\mu$, a Gaussian target has the form
\begin{equation}
 H(\mathbf u)=\exp\left(-\frac{\|\mathbf u-\boldsymbol\mu\|_2^2}
 {2\sigma^2}\right).
 \label{eq:heatmap}
\end{equation}
Here $\mathbf u$ is a two-dimensional grid index, and $\sigma>0$ determines
the target's spread in grid cells. The target is centred on a discrete grid
location; a separate offset compensates for the continuous centre's
displacement from that location. When several objects share a class,
their Gaussian targets are combined by an elementwise maximum. This
heatmap construction originates in CenterNet and is used in CenterPoint
\cite{yin2021centerpoint,zhou2019objectsaspoints}.

The centre head is trained with a variant of focal loss adapted to Gaussian
targets. For target value $H_{uvc}$ and predicted value $\hat H_{uvc}$ at
grid position $(u,v)$ and class channel $c$, a per-location term is
\begin{equation}
 \ell_{uvc}=
 \begin{cases}
 -(1-\hat H_{uvc})^a\log\hat H_{uvc},&H_{uvc}=1,\\
 -(1-H_{uvc})^b(\hat H_{uvc})^a\log(1-\hat H_{uvc}),&H_{uvc}<1.
 \end{cases}
 \label{eq:center_heatmap_loss}
\end{equation}
The exponents $a,b\ge0$ control downweighting, and the complete heatmap
loss sums these terms with normalisation by the number of target centres.
Locations around a Gaussian peak are downweighted as negatives, rather
than treated as additional positive object centres
\cite{zhou2019objectsaspoints}. This supervision differs from the binary
foreground loss in PointRCNN.

At a detected peak, additional heads regress the centre offset, height,
dimensions and orientation~\cite{yin2021centerpoint}. Orientation can be represented
by $(\sin\psi,\cos\psi)$ and decoded by the two-argument arctangent,
which avoids a scalar discontinuity where an angle wraps around.
During training, the box attributes at annotated centres are supervised
with weighted $L_1$ regression terms alongside the centre-heatmap loss.
Each weight scales the contribution of an attribute, such as size or
orientation; the loss therefore supervises localisation as well as the
presence of a centre~\cite{yin2021centerpoint}. At inference, the predicted
attributes are decoded into metric boxes and redundant boxes are
suppressed.

Several input points can occupy the same voxel. Added points may change
the encoded feature of that cell without creating another occupied cell;
the outcome depends on the encoder and point-selection rule
\cite{zhou2018voxelnet}. This distinction explains how point-based and
voxel-based representations process density differently. It is not a
geometric accuracy metric or a claim that one detector is always more
robust. With the downstream processing paths established, the next section
explains how the additional points are generated.

\FloatBarrier
\section{Point Cloud Upsampling}
\label{sec:upsampling}

Classification and detection produce semantic predictions from an input
representation. Upsampling instead changes that representation: it infers
a denser point set from sparse spatial support. Given
$\mathcal X=\{\mathbf x_i\}_{i=1}^N$, a fixed-ratio upsampling model is
written as
\begin{equation}
 \mathcal Y=\mathcal U_\theta(\mathcal X;r)
 =\{\mathbf y_j\}_{j=1}^{rN},\qquad |\mathcal Y|=rN.
 \label{eq:upsampling_map}
\end{equation}
Here $r>1$ is an integer expansion ratio, $\theta$ denotes the parameters
of a learned model, and $j$ indexes output points. This formulation follows
the sparse-to-dense task used by PU-Net~\cite{yu2018punet}; a geometric
optimisation method need not have learned parameters. The output should
represent the underlying surface, preserve its local structure and avoid
excessive clustering. Cardinality alone cannot establish any of those
properties.

It is not presupposed that $\mathcal X\subseteq\mathcal Y$. A coordinate
regression network can predict the whole output anew, while a residual
network can move every repeated input anchor~\cite{yu2018punet,kim2023puedgeformer}.
Thus even an architecture that begins by repeating input coordinates need
not retain the original measurements exactly. Whether observations are
explicitly kept is a separate property of an output construction.

The methods below illustrate five ways of addressing the ambiguity of
sparse-to-dense reconstruction: geometric resampling, hierarchical feature
expansion, graph-based expansion, graph attention and conditional diffusion.
Related developments include adversarial distribution learning in PU-GAN,
progressive patch refinement, and edge-supervised consolidation
\cite{li2019pugan,wang2019mpu,yu2018ecnet}. The following descriptions
separate the published method principles from the wrappers and downstream
training choices introduced later in the thesis.

\subsection{Edge-Aware Resampling (EAR)}
\label{subsec:ear}

EAR is an optimisation-based resampling method developed by Huang et al.
to consolidate noisy point sets while preserving sharp features
\cite{huang2013ear}. The central difficulty is that normals estimated close
to a sharp edge can mix information from different surface sides. EAR
therefore first resamples away from uncertain edge regions to obtain
reliable oriented points. It then inserts points progressively towards
the gaps near edges. This two-stage sequence is the reason for the term
edge-aware resampling; it is more specific than simply smoothing a cloud
and adding points.

Let $\mathbf n_i$ be a unit normal associated with point $\mathbf p_i$.
EAR combines a spatial weight and a normal-similarity weight
\cite{huang2013ear}, giving
\begin{equation}
 \kappa_p(d)=\exp(-d^2/\sigma_p^2),\qquad
 \kappa_n(\mathbf n_i,\mathbf n_j)
 =\exp\left[-\left(\frac{1-\mathbf n_i^{\mathsf T}\mathbf n_j}
 {1-\cos\sigma_n}\right)^2\right].
 \label{eq:ear_kernels}
\end{equation}
Here $d=\|\mathbf p_i-\mathbf p_j\|_2$ is spatial distance,
$\sigma_p>0$ sets its neighbourhood scale, and $0<\sigma_n<\pi$ is an angular
bandwidth in radians. The dot product of unit normals is the cosine of their angular
difference. Nearby points with compatible normals receive high weight;
disagreeing normals are downweighted even when the points are close.
Using $w_{ij}=\kappa_p(\|\mathbf p_i-\mathbf p_j\|_2)
\kappa_n(\mathbf n_i,\mathbf n_j)$, the weighted normal update is
\begin{equation}
 \bar{\mathbf n}_i=
 \frac{\sum_{j\in\mathcal N(i)}w_{ij}\mathbf n_j}
 {\sum_{j\in\mathcal N(i)}w_{ij}}.
 \label{eq:ear_normal_filter}
\end{equation}
The set $\mathcal N(i)$ contains spatial neighbours and $\bar{\mathbf n}_i$
is their weighted mean direction; a nonzero vector can be normalised when
a unit direction is required. This is the bilateral normal-smoothing step,
not the entire EAR algorithm. EAR alternates normal processing with locally
anisotropic projection to move samples away from ambiguous edges, because
normal smoothing alone can retain incorrect assignments
\cite{huang2013ear}.

For a new sample, EAR chooses a base position between existing points,
selects a projection direction compatible with surrounding normals, and
computes a projection distance towards the latent surface. Positions and
normals influence all three decisions. Repeating this process fills sparse
regions and approaches sharp features from better-supported regions
\cite{huang2013ear}. The method does not require a trained generator, but
its geometric neighbourhoods and normals must be meaningful. Sparse or
mixed-surface neighbourhoods therefore warrant inspection when the method
is used outside the surface-scan setting of its original evaluation.
Learned methods address the reconstruction ambiguity differently, through
features estimated from training examples.

\subsection{PU-Net}
\label{subsec:punet}

PU-Net is an early end-to-end network dedicated to point cloud upsampling.
It consists of multi-level feature embedding, feature expansion and
coordinate reconstruction~\cite{yu2018punet}. Its encoder adapts the
hierarchical feature extraction of PointNet++ to local surface patches.
Because deeper levels contain fewer sampled locations, their features are
interpolated back to the original points before aggregation.

For a point $\mathbf x$ and its three nearest feature locations
$\mathbf x_j^{(\ell)}$ at level $\ell$, PU-Net uses inverse-distance
interpolation~\cite{yu2018punet}, expressed as
\begin{equation}
 \tilde{\mathbf f}^{(\ell)}(\mathbf x)=
 \frac{\sum_{j=1}^3 w_j(\mathbf x)\mathbf f^{(\ell)}(\mathbf x_j^{(\ell)})}
 {\sum_{j=1}^3w_j(\mathbf x)},\qquad
 w_j(\mathbf x)=\frac1{\|\mathbf x-\mathbf x_j^{(\ell)}\|_2}.
 \label{eq:punet_interpolation}
\end{equation}
Here $\mathbf f^{(\ell)}$ is the feature field on the sampled level.
The expression applies at positive distances; an exactly coincident
location takes the coincident feature rather than dividing by zero.
Features from the different levels are projected to compatible channel
dimensions and concatenated, giving
$\mathbf F\in\mathbb R^{N\times F}$. The number $F$ denotes the resulting
feature width. Local detail and broader patch context are therefore made
available to the expansion stage at every original point.

PU-Net uses $r$ branches of independent pointwise convolutions, followed
by a reshape~\cite{yu2018punet}, producing
\begin{equation}
 \mathbf F'=\operatorname{RS}\left[
 C_1^2(C_1^1(\mathbf F)),\ldots,C_r^2(C_r^1(\mathbf F))\right]
 \in\mathbb R^{rN\times F_2}.
 \label{eq:punet_expansion}
\end{equation}
Each $C_m^1,C_m^2$ is a learned $1\times1$ convolution in branch $m$.
The bracket concatenates channels into an $N\times rF_2$ tensor;
$\operatorname{RS}$ reshapes it to $rN$ rows of width $F_2$.
Separate branch weights reduce the correlation between the new features.
A shared coordinate head then maps each expanded row to three coordinates,
$\mathbf y_{i,m}=\psi_\theta(\mathbf f'_{i,m})$, with $i$ indexing
an input feature and $m=1,\ldots,r$ its expanded features.

\begin{samepage}
PU-Net uses Earth Mover's Distance (EMD) for reconstruction.
Let $\mathcal Y$ be the predicted point set and $\mathcal R$ its dense
reference, with $|\mathcal Y|=|\mathcal R|$. The reconstruction loss is
\begin{equation}
 \mathcal L_\mathrm{rec}
 =\min_{\substack{\varphi:\mathcal Y\to\mathcal R\\
                  \varphi\ \mathrm{bijective}}}
 \sum_{\mathbf y\in\mathcal Y}\|\mathbf y-\varphi(\mathbf y)\|_2.
 \label{eq:emd}
\end{equation}
\end{samepage}
Here $\mathbf y$ denotes an output point and $\varphi(\mathbf y)$ its
assigned reference point. The bijective map $\varphi$ pairs every output
with exactly one reference point,
unlike nearest-neighbour Chamfer matching, which permits many outputs to
share one reference~\cite{yu2018punet,fan2017pointsetgen}. Equal cardinality
is required by this discrete bijection. Dividing the sum by $|\mathcal Y|$
would give a mean variant with different numerical scaling.

Reconstruction alone can leave clustered outputs. PU-Net therefore adds
the local repulsion term
\begin{equation}
 \mathcal L_\mathrm{rep}=\sum_{i=1}^{rN}\sum_{j\in\mathcal N_k(i)}
 -d_{ij}\exp(-d_{ij}^2/h^2),
 \qquad d_{ij}=\|\mathbf y_i-\mathbf y_j\|_2.
 \label{eq:repulsion}
\end{equation}
Here $\mathcal N_k(i)$ contains the $k$ nearest other output points and
$h>0$ controls the decay scale. The negative distance favours separation
of very close neighbours, while the exponential limits the range over
which the interaction is effective. The published loss combines
reconstruction, weighted repulsion and parameter regularisation
\cite{yu2018punet}. These terms are expressed in the coordinates used for
training. Patch normalisation and its inverse therefore affect the scale
of both reconstructed geometry and local separation. The next method
changes the expansion operation itself by incorporating graph neighbours.

\subsection{PU-GCN}
\label{subsec:pugcn}

PU-GCN proposes an Inception DenseGCN feature extractor and the NodeShuffle
upsampling module~\cite{qian2021pugcn}. Graph convolution represents local
relations between points explicitly. For features $\mathbf h_i$ and
neighbour indices $j\in\mathcal N_k(i)$, the EdgeConv construction used
in this graph-convolution family is
\begin{equation}
 \mathbf e_{ij}=\phi_\theta([\mathbf h_i,\mathbf h_j-\mathbf h_i]),
 \qquad\mathbf h'_i=\max_{j\in\mathcal N_k(i)}\mathbf e_{ij}.
 \label{eq:edge_feature}
\end{equation}
Here $\mathcal N_k(i)$ is the set of indices of the $k$ neighbours of
node $i$; $\mathbf h_i$ and $\mathbf h_j$ are the centre and neighbour
features. The learned shared mapping $\phi_\theta$ transforms their
concatenation into the edge feature $\mathbf e_{ij}$. Channelwise max
aggregation produces the updated node feature $\mathbf h'_i$
\cite{wang2019dgcnn}.
The centre feature preserves information about the point itself, while
the feature difference describes its relationship to a neighbour.
Neighbourhoods may be constructed in a coordinate or feature space; the
space and graph rule are part of the architecture specification.

Inception DenseGCN combines densely connected graph layers with branches
having different receptive fields. Dilated neighbourhoods let the branches
reach different spatial extents without simply duplicating one local
operation. Their features are fused with broader pooled information
\cite{qian2021pugcn}. Dense connections retain intermediate features
so that later processing can use both earlier local detail and deeper
representations. Figure~\ref{fig:pugcn-principle} shows how feature
extraction leads to expansion and coordinate reconstruction.
\thesisfigure{pugcn_principle}{PU-GCN principle}{Graph feature extraction,
NodeShuffle and coordinate reconstruction in PU-GCN, redrawn from the
published method~\cite{qian2021pugcn}.}{fig:pugcn-principle}

The graph layers first embed coordinates and then combine features from
different receptive fields. The block outputs are combined through
residual connections and passed through a pointwise bottleneck before
upsampling~\cite{qian2021pugcn}. Denote the resulting feature of point $i$
by $\mathbf f_i^{\mathrm{multi}}$. Stacking these features as rows gives
$\mathbf F\in\mathbb R^{N\times F}$, where the scalar $F$ is their channel
width. Thus $\mathbf F$ is the processed output of feature extraction,
not an independently supplied input.

NodeShuffle applies a graph convolution to $\mathbf F$ that expands each
node's channel width from $F$ to $rF$,

\begin{equation}
 \mathbf F_\mathrm{exp}=\operatorname{GCN}_\theta(\mathbf F)
 \in\mathbb R^{N\times rF}.
 \label{eq:nodeshuffle_expand}
\end{equation}
The operator $\operatorname{GCN}_\theta$ denotes the graph-convolution
layer with learned parameters $\theta$, and $\mathbf F_\mathrm{exp}$
contains its expanded features. Periodic shuffling then rearranges groups
of $F$ channels into additional rows,
\begin{equation}
 \mathbf F_\mathrm{shuffle}=\operatorname{PS}(\mathbf F_\mathrm{exp})
 \in\mathbb R^{rN\times F}.
 \label{eq:nodeshuffle_ps}
\end{equation}
The matrix $\mathbf F_\mathrm{shuffle}$ contains one feature row per
prospective output point. The operator $\operatorname{PS}$ is a
deterministic rearrangement, not a learned coordinate displacement. The preceding graph convolution learns the new
features; the following coordinate head converts the shuffled features
into dense point positions~\cite{qian2021pugcn}. PU-GCN trains this output
with a squared Chamfer reconstruction loss, whose distance convention is
distinguished from the unsquared metric in Section~\ref{sec:metrics}.

The construction makes neighbourhood selection relevant to transfer.
For fixed $k$, a sparse region can have a much larger neighbourhood radius
than a dense one. A graph that crosses unrelated surfaces presents
different relations to the learned mapping. The graph construction therefore determines which local relationships
reach the learned feature extractor.

Reconstruction training and downstream detector retraining have different
targets. PU-GCN learns dense coordinates using a geometric reconstruction
loss~\cite{qian2021pugcn}. PointRCNN and CenterPoint learn object
classification and box localisation through their detection losses
\cite{shi2019pointrcnn,yin2021centerpoint}. With a fixed upsampler, its
outputs can be supplied as detector-training inputs while the detection
loss updates only the detector parameters. The saved point-cloud files are not regenerated by that optimisation.

The detector-side retraining studied in this thesis follows this
separation: PointRCNN and CenterPoint are adapted to PU-GCN inputs,
while the PU-GCN upsampling network retains its fixed checkpoint.
Section~\ref{sec:fc-4-11} specifies the completed training conditions;
Figure~\ref{fig:adaptation-design} identifies which network is updated.

Attention-based upsampling changes feature generation itself by combining
local relationships with the wider patch context.

\subsection{PU-EdgeFormer}
\label{subsec:puedgeformer}

PU-EdgeFormer combines local graph features with multi-head
self-attention~\cite{kim2023puedgeformer}. Its encoder stacks EdgeFormer
units, followed by feature extension and coordinate reconstruction. To
explain the modification, standard scaled dot-product attention first
forms queries, keys and values from a feature matrix
$\mathbf H\in\mathbb R^{N\times F}$ through learned projections,
\begin{equation}
 \mathbf Q=\mathbf H\mathbf W_Q,\quad
 \mathbf K=\mathbf H\mathbf W_K,\quad
 \mathbf V=\mathbf H\mathbf W_V,
 \label{eq:qkv}
\end{equation}
\begin{equation}
 \operatorname{Attn}(\mathbf Q,\mathbf K,\mathbf V)
 =\operatorname{softmax}\left(\frac{\mathbf Q\mathbf K^{\mathsf T}}
 {\sqrt{d_k}}\right)\mathbf V.
 \label{eq:attention}
\end{equation}
The learned matrices $\mathbf W_Q,\mathbf W_K$ produce $d_k$-dimensional
query and key vectors; $\mathbf W_V$ produces value vectors. The dot
product scores relationships between locations. Rowwise softmax turns
each row into nonnegative weights summing to one, and multiplication by
$\mathbf V$ gives a weighted feature combination. Scaling by
$\sqrt{d_k}$ controls the magnitude of the dot products. Multiple heads
perform these operations in separate learned subspaces, concatenate their
outputs and apply an output projection~\cite{vaswani2017attention}.

The defining change in PU-EdgeFormer is that EdgeConv replaces the linear
query, key and value transformations. In an architectural shorthand,
\begin{equation}
 \mathbf Q=\operatorname{EdgeConv}_Q(\mathbf H),\quad
 \mathbf K=\operatorname{EdgeConv}_K(\mathbf H),\quad
 \mathbf V=\operatorname{EdgeConv}_V(\mathbf H).
 \label{eq:edgeformer_qkv}
\end{equation}
Each branch groups nearest neighbours, convolves the grouped features and
max-pools them. The resulting query, key and value features therefore
contain local relationships before attention combines information across
the patch~\cite{kim2023puedgeformer}. 

The feature-extension stage shuffles an $N\times F$ tensor into
$rN\times(F/r)$ features and applies MLPs, giving a dense feature matrix
$\mathbf F'$. If $\tilde{\mathbf X}\in\mathbb R^{rN\times3}$ repeats
each input coordinate $r$ times, the coordinate reconstruction is
\begin{equation}
 \mathbf Y=\tilde{\mathbf X}+\Phi_\theta(\mathbf F').
 \label{eq:edgeformer_reconstruction}
\end{equation}
The mapping $\Phi_\theta$ predicts three-dimensional offsets, and the rows
of $\mathbf Y$ form the output set~\cite{kim2023puedgeformer}. Repeated
anchors do not imply repeated outputs because their predicted offsets can
differ. Equally, residual reconstruction does not guarantee exact retention
of the measured coordinates. The method trains with squared Chamfer
distance. Attention expands the available context within the input patch,
but cannot establish whether that patch corresponds to a coherent observed
surface. Conditional diffusion takes a different route by generating the
output through successive denoising steps.

\subsection{PDANS}
\label{subsec:pdans}

PDANS combines a conditional diffusion model with Adaptive Noise
Suppression (ANS) and a TreeTrans decoder~\cite{zhang2025pdans}. Its starting
point is the denoising diffusion probabilistic model (DDPM), which learns
to reverse a gradual Gaussian noising process~\cite{ho2020ddpm}. The
sparse point cloud conditions this reversal so that the generated dense
points correspond to the supplied geometry.

Let $\mathbf Z_0\in\mathbb R^{3rN}$ be the vectorised dense point
coordinates, and let $t=1,\ldots,T$ index noising steps. A forward step is
\begin{equation}
 q(\mathbf Z_t\mid\mathbf Z_{t-1})=
 \mathcal N(\sqrt{1-\beta_t}\mathbf Z_{t-1},\beta_t\mathbf I),
 \label{eq:diffusion_forward}
\end{equation}
where $\beta_t\in(0,1)$ is the variance increment, $\mathbf I$ is the
identity in $3rN$ dimensions, and $\mathcal N(\boldsymbol\mu,\boldsymbol\Sigma)$
denotes a Gaussian with mean $\boldsymbol\mu$ and covariance
$\boldsymbol\Sigma$. With $\alpha_t=1-\beta_t$ and
$\bar\alpha_t=\prod_{s=1}^t\alpha_s$, the state can be sampled directly as
\begin{equation}
 \mathbf Z_t=\sqrt{\bar\alpha_t}\mathbf Z_0+
 \sqrt{1-\bar\alpha_t}\boldsymbol\epsilon,
 \qquad\boldsymbol\epsilon\sim\mathcal N(\mathbf0,\mathbf I).
 \label{eq:diffusion_direct}
\end{equation}
The product $\bar\alpha_t$ is the remaining signal fraction after $t$
steps; the product index $s$ runs over the forward steps up to $t$,
and $\boldsymbol\epsilon$ is standard Gaussian noise. These are the
DDPM forward equations also used to introduce PDANS
\cite{zhang2025pdans,ho2020ddpm}.

If $\mathbf c$ denotes the conditioning representation of the sparse
input and requested expansion, a learned reverse step is
\begin{equation}
 p_\theta(\mathbf Z_{t-1}\mid\mathbf Z_t,\mathbf c)
 =\mathcal N(\boldsymbol\mu_\theta(\mathbf Z_t,t,\mathbf c),
 \sigma_t^2\mathbf I).
 \label{eq:diffusion_reverse}
\end{equation}
Here $\boldsymbol\mu_\theta$ is the learned reverse mean and $\sigma_t^2$
the reverse variance. The conditioning vector $\mathbf c$ encodes the
sparse input and requested expansion. In the noise-prediction
parameterisation, $\boldsymbol\epsilon_\theta$ is the learned function
that estimates the noise present at step $t$, given the noisy coordinates
and their conditioning. It is trained through the mean-squared objective
\begin{equation}
 \mathcal L_\mathrm{diff}=\mathbb E_{\mathbf Z_0,t,\boldsymbol\epsilon}
 \left[\|\boldsymbol\epsilon-
 \boldsymbol\epsilon_\theta(\mathbf Z_t,t,\mathbf c)\|_2^2\right].
 \label{eq:diffusion_loss}
\end{equation}
The expectation averages over training shapes, sampled times and sampled
noise; $\mathbf Z_t$ is constructed by Equation~\eqref{eq:diffusion_direct}.
At inference the process starts from noise and repeatedly estimates a
less noisy state. Unlike a direct coordinate head, generation therefore
depends on a sequence of denoising updates. The squared noise-prediction
loss supervises the denoising process. Geometric metrics instead assess
the final generated coordinates, as described in Section~\ref{sec:metrics}
\cite{zhang2025pdans,ho2020ddpm}.

The ANS module constructs neighbourhoods for the input points and uses attention
to calculate weights for refining coordinates and features
\cite{zhang2025pdans}. For centre $\mathbf x_i$ and neighbour
$\mathbf x_{ij}$, write the coordinate-based attention score as
\begin{equation}
 a_{ij}=\frac{\exp(s_{ij})}{\sum_{m\in\mathcal N_k(i)}\exp(s_{im})},
 \qquad s_{ij}=\frac{q_\theta(\mathbf x_i)^{\mathsf T}
 k_\theta(\mathbf x_{ij}-\mathbf x_i)}{\sqrt{d_q}}.
 \label{eq:pdans_attention}
\end{equation}
Here $q_\theta$ and $k_\theta$ are learned coordinate transformations,
$d_q$ is their channel dimension, and $\mathcal N_k(i)$ indexes the $k$
neighbours. The query uses the centre coordinates and the key uses
relative neighbour coordinates. With a learned linear feature transform
$\gamma_\theta$, the weighted neighbourhood representation is
\begin{equation}
 \mathbf w_i=\sum_{j\in\mathcal N_k(i)}a_{ij}\gamma_\theta(\mathbf f_{ij}).
 \label{eq:pdans_weights}
\end{equation}
The vector $\mathbf f_{ij}$ is the neighbour feature. ANS uses learned
activations of $\mathbf w_i$ to rescale both the centre coordinates and
its features, then fuses the refined information. This is a learned
suppression mechanism; the attention weights are not independently
validated probabilities that a point is noise~\cite{zhang2025pdans}.

TreeTrans propagates features from coarser tree nodes to finer descendants
using transformations of the current node and information from ancestor
nodes. Local and global attention then combine the tree features with
corresponding encoder features. The conditional and diffusion networks
use these representations to guide dense-point generation
\cite{zhang2025pdans}. ANS addresses unreliable input support, while
TreeTrans improves how features at different levels enter the decoder.
The resulting surface quality must still be checked against downstream
predictions; a denoising objective does not directly optimise class labels
or detection boxes. This brings the discussion to evaluation: the metrics
must distinguish geometric fidelity, point distribution and task success.

\FloatBarrier
\section{Evaluation Metrics}
\label{sec:metrics}

The methods above produce point coordinates, whereas the downstream
models produce labels or boxes. These outputs answer different questions
and require different metrics. Geometric evaluation compares a point set
with a reference surface or reference samples. Classification and detection
evaluation compare predictions with semantic annotations
\cite{qian2021pugcn,shi2019pointrcnn,qi2017pointnet}. Their definitions are
given here; the implemented aggregation rules and experimental reference
construction are specified in Chapter~\ref{chap:setup}.

\subsection{Geometric Metrics}
\label{subsec:metrics_geometric}

Let $\mathcal Y$ be a nonempty output point set and $\mathcal R$ a nonempty
reference set. Define the distance from a point $\mathbf y$ to its nearest
reference point as
\begin{equation}
 d(\mathbf y,\mathcal R)=\min_{\mathbf r\in\mathcal R}
 \|\mathbf y-\mathbf r\|_2.
 \label{eq:nearest_reference}
\end{equation}
This nearest-neighbour construction is used in Chamfer and Hausdorff
comparisons~\cite{qian2021pugcn,fan2017pointsetgen}. It requires a common
coordinate frame. If the coordinates are measured in metres, an unsquared
distance is in metres; after unit-sphere normalisation it is dimensionless.
Squaring the distance changes its units and its sensitivity to large errors.

A symmetric averaged, unsquared Chamfer distance is
\begin{equation}
 \begin{aligned}
 \operatorname{CD}(\mathcal Y,\mathcal R)
 ={}&\frac1{|\mathcal Y|}\sum_{\mathbf y\in\mathcal Y}
 \min_{\mathbf r\in\mathcal R}\|\mathbf y-\mathbf r\|_2\\
 &+\frac1{|\mathcal R|}\sum_{\mathbf r\in\mathcal R}
 \min_{\mathbf y\in\mathcal Y}\|\mathbf r-\mathbf y\|_2.
 \end{aligned}
 \label{eq:fc-2-10}
\end{equation}
The cardinality symbols $|\mathcal Y|$ and $|\mathcal R|$ denote the
numbers of points. The first term measures how far outputs lie from the
reference; the second measures how well the output covers the reference.
Both directions are necessary: points concentrated on one small part of
a surface can have a small first term while leaving much of the reference
uncovered. This bidirectional construction follows the nearest-neighbour
matching used for point-set reconstruction
\cite{fan2017pointsetgen,qian2021pugcn}. Because the terms are averages,
a small number of severe deviations can be hidden by many close matches.

Distance conventions must be distinguished explicitly. Fan et al. define
the sum of squared nearest-neighbour distances, and PU-GCN uses averaged
squared Euclidean distances in its Chamfer loss
\cite{qian2021pugcn,fan2017pointsetgen}. Equation~\eqref{eq:fc-2-10} defines
an unsquared, averaged variant, so it must not be cited as numerically
identical to those losses. Replacing each norm by its square produces the
averaged squared variant; multiplying a reported value by a display factor
does not convert one variant into the other.

The symmetric Hausdorff distance takes the worst nearest-neighbour error
in either direction~\cite{huttenlocher1993hausdorff}, with
\begin{equation}
 \begin{aligned}
 h(\mathcal Y,\mathcal R)&=
 \max_{\mathbf y\in\mathcal Y}\min_{\mathbf r\in\mathcal R}
 \|\mathbf y-\mathbf r\|_2,\\
 \operatorname{HD}(\mathcal Y,\mathcal R)&=
 \max\{h(\mathcal Y,\mathcal R),h(\mathcal R,\mathcal Y)\}.
 \end{aligned}
 \label{eq:fc-2-11}
\end{equation}
Here $h$ is the directed distance; reversing its arguments exchanges output
deviation and reference coverage. HD highlights isolated outliers and
uncovered regions that may be diluted in CD. It is correspondingly sensitive
to the worst reference error. PU-GCN reports HD alongside CD and P2F to
characterise different geometric properties~\cite{qian2021pugcn}.

If a continuous reference mesh surface $S$ is available, a mean unsquared
point-to-surface distance is
\begin{equation}
 \operatorname{P2F}(\mathcal Y,S)=\frac1{|\mathcal Y|}
 \sum_{\mathbf y\in\mathcal Y}\min_{\mathbf s\in S}\|\mathbf y-\mathbf s\|_2.
 \label{eq:p2f}
\end{equation}
Here $\mathbf s$ ranges over the surface, including face interiors, rather
than only its vertices. Such mesh-based distances are used in point-cloud
consolidation and upsampling evaluation~\cite{qian2021pugcn,yu2018ecnet}.
They reduce the influence of how densely a discrete reference is sampled.
However, P2F is one-directional: a few points on the surface can obtain a
small value without covering the whole surface. The mesh and output must
share a coordinate frame and scale, and the mesh must represent the intended
target geometry.

Geometric fidelity does not fully describe distribution. PU-Net uses local
repulsion to discourage clustered predictions, while PU-GAN includes a
separate uniformity objective~\cite{yu2018punet,li2019pugan}.
A basic local spacing quantity is
\begin{equation}
 d_i^\mathrm{NN}=\min_{j\ne i}\|\mathbf y_i-\mathbf y_j\|_2.
 \label{eq:nnd}
\end{equation}
Here $i$ and $j$ index distinct stored rows, so repeated coordinates can
produce zero spacing. Neighbour distances underlie the local repulsion
and uniformity constructions in PU-Net and PU-GAN
\cite{yu2018punet,li2019pugan}. A nearest-neighbour distribution, a disk-count
uniformity statistic and a repulsion loss describe related but different
properties. Their neighbourhood definition, scale and aggregation must be
given before their numbers are compared. A metric abbreviated as a
uniformity coefficient should not silently be replaced by another such
statistic.

CD and HD are mathematically defined for nonempty sets of unequal
cardinality; the discrete EMD bijection in Equation~\eqref{eq:emd} is not.
Equal point counts in a method comparison are an experimental control,
not a condition needed to evaluate a nearest-neighbour minimum. Changing
the sample count or reference density can change coverage and matching
distances, so both must be reported. There is no general theorem that
adding arbitrary points improves the symmetric CD: additional outliers
can worsen its output-to-reference term. The implemented equal-count
protocol, reference construction and uniformity statistic are defined in
Chapter~\ref{chap:setup}.

\subsection{Classification Metrics}
\label{subsec:metrics_classification}

Point-cloud classification studies report overall accuracy (OA) and mean
class accuracy (mAcc)~\cite{qi2017pointnet}. For $M$ evaluated samples
with true labels $y_m$ and predicted labels $\hat y_m$, OA is
\begin{equation}
 \operatorname{OA}=\frac1M\sum_{m=1}^M\mathbb I[\hat y_m=y_m].
 \label{eq:fc-2-12}
\end{equation}
The indicator $\mathbb I[\cdot]$ equals one when the condition holds and
zero otherwise; $m$ indexes an object, not an individual point. OA weights
each evaluated object equally. For class $c$ containing $M_c>0$ objects,
\begin{equation}
 \operatorname{Acc}_c=\frac1{M_c}\sum_{m:y_m=c}\mathbb I[\hat y_m=c],
 \qquad
 \operatorname{mAcc}=\frac1C\sum_{c=1}^C\operatorname{Acc}_c.
 \label{eq:macc}
\end{equation}
Here $C$ is the number of evaluated classes. mAcc gives each class equal
weight regardless of its sample count, whereas OA gives more aggregate
weight to larger classes. They can therefore change by different amounts
on an imbalanced evaluation set~\cite{qi2017pointnet}.

These equations define sample-level aggregation. Averaging batch
accuracies without weighting by batch size gives a different quantity when
batches differ in size. Similarly, averaging only the class entries present
in each batch can differ from pooling all class predictions first. The
actual evaluator must be read to determine which aggregation produced a
reported result; the implementation used here is stated explicitly in
Chapter~\ref{chap:setup}. This distinction prevents a textbook definition
from being substituted for an executed metric.

\subsection{Detection Metrics}
\label{subsec:metrics_detection}

Detection evaluates class correctness together with localisation. For
predicted and ground-truth boxes $B_\mathrm{pred}$ and $B_\mathrm{gt}$, the
three-dimensional intersection over union is
\begin{equation}
 \operatorname{IoU}_{3D}=
 \frac{\operatorname{Vol}(B_\mathrm{pred}\cap B_\mathrm{gt})}
 {\operatorname{Vol}(B_\mathrm{pred}\cup B_\mathrm{gt})}.
 \label{eq:iou3d}
\end{equation}
The operators $\cap$ and $\cup$ denote intersection and union, and
$\operatorname{Vol}$ is volume. For nondegenerate boxes IoU lies between
zero and one. Replacing the volumes by the areas of horizontal box
footprints gives BEV IoU. A BEV overlap can remain high when vertical
localisation is inaccurate, so BEV and full 3D evaluation should not be
interchanged~\cite{shi2019pointrcnn,simonelli2019disentangling}.

Predictions are ranked by confidence and matched to eligible ground-truth
instances according to class, overlap and benchmark rules. A matched
prediction is a true positive (TP); an unmatched eligible prediction is a
false positive (FP), and a missed eligible object is a false negative (FN).
At a given score threshold,
\begin{equation}
 \operatorname{Precision}=\frac{TP}{TP+FP},
 \qquad
 \operatorname{Recall}=\frac{TP}{TP+FN}.
 \label{eq:precision_recall}
\end{equation}
The denominators count retained eligible predictions and eligible annotated
objects, respectively. Duplicate detections cannot all match the same
object as true positives; ignored annotations and empty denominators must
be handled according to the benchmark implementation. Sweeping the score
threshold produces a precision--recall curve~\cite{everingham2010voc}.

Interpolated precision at recall $\nu$ is the upper envelope
\begin{equation}
 p_\mathrm{interp}(\nu)=\max_{\tilde\nu\ge\nu}p(\tilde\nu),
 \label{eq:interp_precision}
\end{equation}
where $p(\tilde\nu)$ is the precision achieved at an attained recall
$\tilde\nu$. The maximum removes local upward fluctuations in precision
\cite{everingham2010voc}. Average precision summarises the interpolated precision across the
recall positions specified by an evaluation protocol
\cite{everingham2010voc,simonelli2019disentangling}. Different sampling
conventions need not yield the same numerical AP. The implemented
40-position convention, class thresholds and difficulty rules are
specified in Section~\ref{sec:fc-4-10-3}.

An individual scene image can illustrate a missed or recovered detection,
but it does not determine AP for a complete evaluation set. Likewise,
changing confidence ordering can affect AP even if some selected boxes
appear similar. Qualitative examples and aggregate precision--recall
evaluation therefore play complementary roles.

The geometric and task metrics quantify different outputs. Their
relationship is examined through the comparison design in
Chapter~\ref{chap:concept}; the datasets and executed evaluation settings
are specified in Chapter~\ref{chap:setup}.
\FloatBarrier
